\documentclass[12pt,a4paper]{article}
\usepackage[top=2.5cm,bottom=2.5cm,left=2.5cm,right=2.5cm,headheight=15pt]{geometry}
\usepackage{amsmath,amssymb}
\usepackage{tikz}
\usetikzlibrary{shapes,arrows,arrows.meta,positioning,calc,backgrounds,fit,matrix}
\usepackage{booktabs,array,longtable,tabularx,multirow}
\usepackage{xcolor}
\usepackage{enumitem}
\usepackage{fancyhdr}
\usepackage{titlesec}
\usepackage{mdframed}
\usepackage{tcolorbox}
\tcbuselibrary{skins,breakable}
\usepackage{hyperref}

%% ---- Colors ----
\definecolor{folblue}{RGB}{0,82,165}
\definecolor{lightblue}{RGB}{220,235,255}
\definecolor{darkorange}{RGB}{180,70,0}
\definecolor{lightorange}{RGB}{255,235,210}
\definecolor{darkgreen}{RGB}{20,110,20}
\definecolor{lightgreen}{RGB}{200,240,205}
\definecolor{darkred}{RGB}{150,20,20}
\definecolor{lightred}{RGB}{255,210,210}
\definecolor{lightgray}{RGB}{242,242,242}

%% ---- tcolorbox styles ----
\newtcolorbox{qbox}[2][]{enhanced,breakable,
  colback=lightorange,colframe=darkorange,arc=4pt,
  title={\textbf{#2}},fonttitle=\bfseries\color{white},
  attach boxed title to top left={yshift=-2mm,xshift=4mm},
  boxed title style={colback=darkorange,arc=3pt},#1}

\newtcolorbox{solbox}[2][]{enhanced,breakable,
  colback=lightblue,colframe=folblue,arc=4pt,
  title={\textbf{#2}},fonttitle=\bfseries\color{white},
  attach boxed title to top left={yshift=-2mm,xshift=4mm},
  boxed title style={colback=folblue,arc=3pt},#1}

\newtcolorbox{keybox}[2][]{enhanced,breakable,
  colback=lightgreen,colframe=darkgreen,arc=4pt,
  title={\textbf{#2}},fonttitle=\bfseries\color{white},
  attach boxed title to top left={yshift=-2mm,xshift=4mm},
  boxed title style={colback=darkgreen,arc=3pt},#1}

\newtcolorbox{algobox}[2][]{enhanced,breakable,
  colback=lightgray,colframe=gray!60,arc=4pt,
  title={\textbf{#2}},fonttitle=\bfseries,
  attach boxed title to top left={yshift=-2mm,xshift=4mm},
  boxed title style={colback=gray!60,arc=3pt},#1}

\newtcolorbox{warnbox}[2][]{enhanced,breakable,
  colback=lightred,colframe=darkred,arc=4pt,
  title={\textbf{#2}},fonttitle=\bfseries\color{white},
  attach boxed title to top left={yshift=-2mm,xshift=4mm},
  boxed title style={colback=darkred,arc=3pt},#1}

%% ---- Page style ----
\pagestyle{fancy}\fancyhf{}
\lhead{\color{folblue}\textbf{Blockchain \& Bitcoin --- CSE 717 InfoSec}}
\rhead{\color{folblue}Exam: 17 Jun 2026}
\cfoot{--- \thepage\ ---}
\renewcommand{\headrulewidth}{0.4pt}

%% ---- Section style ----
\titleformat{\section}{\large\bfseries\color{folblue}}{}{0em}{}[\titlerule]
\titleformat{\subsection}{\normalsize\bfseries\color{darkorange}}{}{0em}{}

\begin{document}

\begin{center}
{\LARGE\bfseries\color{folblue} CSE 717 --- Information Security}\\[4pt]
{\large\bfseries Topic 12: Blockchain \& Bitcoin --- Architecture, Hashing/PoW, Bitcoin Mining}\\[2pt]
{\large Complete Past Paper Solutions: 2020--2024}\\[6pt]
{\normalsize Source: Lc\#11A (Blockchain), Lc\#11B (Bitcoin Mining)}\\[2pt]
{\normalsize\color{darkred}\textbf{2/5 years confirmed (2021, 2024), 3--8.75 marks. NEW dedicated 2026 materials --- expect strong emphasis.}}
\end{center}

\tableofcontents

%% ==========================================
\section*{Quick Reference --- Blockchain \& Bitcoin (CHEAT SHEET)}
%% ==========================================

\begin{algobox}{Block structure (Lc\#11A p2)}
Every block has a header containing:
\begin{enumerate}[noitemsep]
  \item \textbf{Version} --- version of the block validation rules.
  \item \textbf{Previous Block Header Hash} --- reference to this block's parent --- \textbf{this is what chains the blocks together}.
  \item \textbf{Merkle Root Hash} --- a cryptographic hash of ALL transactions in this block (hash of hashes, combined pairwise up a tree).
  \item \textbf{Time} --- approximate creation time.
  \item \textbf{nBits} --- the current Proof-of-Work difficulty target.
  \item \textbf{Nonce ("number used once")} --- a random counter miners vary to find a valid block hash.
\end{enumerate}
Below the header: the \textbf{coinbase transaction} (miner's reward) + the list of \textbf{transactions}.
\end{algobox}

\begin{keybox}{How hashing makes the chain secure (Lc\#11A p3--4)}
\begin{itemize}[noitemsep]
  \item Each block's hash is a unique "fingerprint" of its contents. ANY change inside a block changes its hash.
  \item Block $n$ stores the hash of block $n-1$ as its \textbf{Previous Hash} field.
  \item If an attacker edits block 2's data, Block 2's hash changes $\Rightarrow$ Block 3's stored "Previous Hash" no longer matches Block 2's new hash $\Rightarrow$ the tamper is immediately detectable.
  \item To successfully tamper, an attacker would have to recompute the hashes of \textbf{every subsequent block} --- and (in a real network) outrun the honest nodes' Proof-of-Work, which is computationally infeasible.
\end{itemize}
\end{keybox}

\begin{algobox}{Proof of Work (PoW) (Lc\#11A p4)}
\begin{itemize}[noitemsep]
  \item A PoW is a \textbf{computational puzzle}: find a \textbf{nonce} such that $\text{hash}(\text{block header} \,\|\, \text{nonce})$ produces a hash below a target value (e.g.\ starts with a required number of leading zeros).
  \item This is hard to solve but \textbf{easy for others to verify} (just hash once and check).
  \item Miners repeatedly try different nonces until one works --- this is "mining."
  \item \textbf{Why it secures the chain:} to alter a past block, an attacker must redo its PoW \emph{and} the PoW of every block after it, faster than the rest of the network combined --- effectively impossible once a block has several confirmations.
\end{itemize}
\end{algobox}

\begin{keybox}{Merkle Tree --- data-integrity check (Lc\#11A p2)}
\begin{itemize}[noitemsep]
  \item Each transaction is hashed individually (leaf nodes).
  \item Pairs of hashes are concatenated and hashed again, repeatedly, up to a single \textbf{Merkle root}.
  \item The Merkle root is stored in the block header.
  \item \textbf{Integrity check:} if ANY transaction in the block changes, the Merkle root changes $\Rightarrow$ mismatch with the header $\Rightarrow$ tampering detected \emph{without} re-checking every transaction individually.
\end{itemize}
\end{keybox}

\begin{algobox}{Distributed P2P Network \& consensus (Lc\#11A p5--6)}
\begin{enumerate}[noitemsep]
  \item Every node in the P2P network holds a full (or partial) copy of the chain.
  \item To add a block: (1) one node will need to update all chains on the network; (2) redo the proof-of-work for each block; (3) take control of more than 50\% of the peer-to-peer network.
  \item \textbf{Consensus} = nodes agree on what blocks are valid; \textbf{decentralization} means no single point of failure.
\end{enumerate}
\textbf{4 pillars of blockchain security (Lc\#11A p7):} Resilience (no single point of failure), Time-stamping, Reliability (certified transactions), Unchangeable transactions (immutability).
\end{algobox}

\begin{keybox}{Types of blockchain (Lc\#11A p8--9)}
\begin{itemize}[noitemsep]
  \item \textbf{Public:} fully decentralized, anyone can join/read/write (e.g.\ Bitcoin, Ethereum).
  \item \textbf{Private:} access restricted, controlled by one organization (e.g.\ Hyperledger, R3 Corda).
  \item \textbf{Hybrid:} combines public security/transparency with private confidentiality (e.g.\ Dragonchain).
\end{itemize}
\end{keybox}

\begin{algobox}{Blockchain vs.\ Shared Database (Lc\#11A p10)}
\begin{center}
\begin{tabular}{@{}p{3.2cm}p{4.8cm}p{4.8cm}@{}}
\toprule
\textbf{Parameter} & \textbf{Blockchain} & \textbf{Shared Database} \\
\midrule
Operations & Insert only & Create/Read/Update/Delete \\
Replication & Full replication on every peer & Master-slave / multi-master \\
Consensus & Most peers agree on outcome & 2-phase commit / Paxos \\
Validation & Global rules enforced network-wide & Only local integrity constraints \\
Disintermediation & Allowed & Not allowed \\
Confidentiality & Fully confidential & Not totally confidential \\
\bottomrule
\end{tabular}
\end{center}
\end{algobox}

\begin{algobox}{Bitcoin essentials (Lc\#11B)}
\begin{itemize}[noitemsep]
  \item \textbf{Bitcoin $\neq$ Blockchain.} Bitcoin is a consensus network + digital currency; blockchain is the underlying technology that records its transactions.
  \item \textbf{Mining} = competing to add the next block by solving the PoW puzzle; the winner broadcasts the block and collects the \textbf{block reward} + transaction fees.
  \item \textbf{Mining difficulty} self-adjusts so that, on average, one block is found every $\sim$10 minutes (Bitcoin), regardless of total network hash power.
  \item \textbf{Mining pools} let many miners combine hash power and split rewards proportionally --- reduces reward variance.
  \item \textbf{Wallet} = software holding your private/public key pair; controls access to funds recorded on-chain (the coins themselves never "move" off-chain).
  \item \textbf{Bitcoin exchange} = where BTC is bought/sold for fiat or other crypto (e.g.\ Binance, Coinbase).
\end{itemize}
\end{algobox}

%% ==========================================
\section{2024 Q3(b) --- Blockchain Security via Hashing + Proof of Work [3 marks]}
%% ==========================================

\begin{qbox}{Question --- 2024 Q3(b) [3]}
A company maintains a blockchain-based distributed ledger for its digital transactions. Each new block contains the hash of the previous block, forming a chain. To add a block, participants must solve a Proof of Work puzzle, making tampering computationally difficult. Now, explain how blockchain ensures security, specifically in terms of hashing and Proof of Work.
\end{qbox}

\begin{solbox}{Answer}
\textbf{(1) Security via hashing (chain-linking):}
\begin{itemize}[noitemsep]
  \item Each block's header stores the \textbf{cryptographic hash of the previous block}, creating a tamper-evident chain.
  \item A block's own hash is computed over all of its contents (header fields + Merkle root of its transactions). Any change to even a single transaction changes the Merkle root, which changes the block's hash.
  \item If an attacker modifies block $k$, block $k$'s hash changes, so the "previous hash" field stored in block $k+1$ no longer matches --- the chain is immediately broken/detectable.
  \item To make the forged chain look valid again, the attacker must recompute the hashes (and headers) of \textbf{every block from $k$ to the latest block}.
\end{itemize}

\textbf{(2) Security via Proof of Work:}
\begin{itemize}[noitemsep]
  \item Each block's hash must satisfy a difficulty target (e.g.\ a required number of leading zero bits); finding such a hash requires brute-forcing a \textbf{nonce} --- computationally expensive but cheap to verify.
  \item Recomputing one block's PoW takes significant time/energy; recomputing the PoW for \emph{all} subsequent blocks (to cover up a tamper) is exponentially harder.
  \item Meanwhile, honest miners keep extending the real chain, so the attacker's forged chain falls further behind and is rejected by the network's "longest valid chain" rule.
\end{itemize}

\textbf{Conclusion:} hashing makes any tampering \emph{detectable} (chain of fingerprints), and Proof of Work makes tampering \emph{computationally infeasible to hide} (every block after the tampered one must be re-mined faster than the honest network) --- together they give the ledger its immutability and integrity guarantees.
\end{solbox}

%% ==========================================
\section{2021 Section B Q3 --- Blockchain/Cryptography Relation, Bitcoin, I/O Script Validation [8.75 marks]}
%% ==========================================

\begin{qbox}{Question --- 2021 Section B Q3(a) [2.75 + 2]}
(i) How is the blockchain related with cryptography? [2.75]\\
(ii) What do you understand about bitcoin? [2]
\end{qbox}

\begin{solbox}{Answer (i) --- Blockchain $\leftrightarrow$ Cryptography [2.75]}
Blockchain relies on cryptography at every layer:
\begin{itemize}[noitemsep]
  \item \textbf{Cryptographic hash functions} link blocks together (each block stores the hash of the previous block, forming the chain) and build the \textbf{Merkle tree} that summarizes all transactions in a block --- this gives \textbf{integrity and tamper-evidence}.
  \item \textbf{Public-key cryptography (asymmetric keys)}: each participant has a key pair. Transactions are \textbf{digitally signed} with the sender's private key and verified with their public key --- this gives \textbf{authentication} (proof of ownership) and \textbf{non-repudiation}.
  \item \textbf{Proof of Work} is itself a cryptographic-hashing puzzle (find a nonce so the hash meets a difficulty target) --- this gives \textbf{consensus/security} against tampering.
\end{itemize}
$\Rightarrow$ In short: hashing secures the chain's structure and data integrity; public-key cryptography (digital signatures) secures ownership and authorization of transactions.
\end{solbox}

\begin{solbox}{Answer (ii) --- What is Bitcoin? [2]}
\begin{itemize}[noitemsep]
  \item Bitcoin is a \textbf{decentralized digital currency / peer-to-peer electronic cash system} with no central bank or single administrator.
  \item Transactions are verified by network nodes via cryptography (digital signatures) and recorded on a public distributed ledger called a \textbf{blockchain}.
  \item New bitcoins are created (and transactions confirmed) through \textbf{mining} --- nodes competing to solve a Proof-of-Work puzzle.
  \item \textbf{Note:} Bitcoin (the currency/network) is not the same as blockchain (the underlying technology) --- bitcoin uses blockchain, but not all blockchains are bitcoin.
\end{itemize}
\end{solbox}

\begin{qbox}{Question --- 2021 Section B Q3(b) [4]}
Discuss the process of input-output scripts for signature validation for bitcoin.
\end{qbox}

\begin{solbox}{Answer --- Input/Output Scripts \& Signature Validation [4]}
Every bitcoin transaction consists of \textbf{inputs} (references to coins being spent) and \textbf{outputs} (new coins being created/assigned to recipients). Spending is controlled by a tiny stack-based scripting system:

\begin{enumerate}[noitemsep]
  \item \textbf{Output script (``locking script'' / scriptPubKey):} attached to each transaction output, it specifies the \textbf{condition} that must be satisfied to spend that output --- typically "a valid signature matching a given public key/address."
  \item \textbf{Input script (``unlocking script'' / scriptSig):} when the recipient later wants to spend that output, they provide an input script containing their \textbf{digital signature} (and public key).
  \item \textbf{Validation process:}
  \begin{enumerate}[noitemsep]
    \item The node concatenates/executes the \textbf{input script} followed by the referenced \textbf{output script}.
    \item The script engine checks: does the supplied public key hash match the one in the output script? Does the supplied signature verify correctly against the transaction data using that public key?
    \item If the combined script evaluates to \textbf{TRUE} (signature valid + conditions met), the output is successfully "unlocked" and can be spent; the new transaction's outputs become locked with new conditions for the next owner.
  \end{enumerate}
  \item This mechanism proves the spender \textbf{owns the private key} corresponding to the coin's address \textbf{without revealing the private key itself} --- the core of bitcoin's ownership/authorization model.
\end{enumerate}
\end{solbox}

%% ==========================================
\section*{2020, 2022, 2023 --- No Blockchain/Bitcoin question found}
%% ==========================================
\begin{warnbox}{2020/2022/2023 papers}
No Blockchain or Bitcoin question was identified in the 2020, 2022, or 2023 past papers (verified against the consolidated past-papers PDF). \textbf{Note:} an earlier tracker entry claimed a 12-mark "2022 Q8" Merkle-tree/ETH-vs-BTC question --- this was checked against the actual 2022 page (which contains only Q1--Q5) and does \emph{not} exist; it has been removed. Blockchain \& Bitcoin is confirmed in 2/5 years (2021, 2024), with brand-new dedicated 2026 lecture materials (Lc\#11A/11B) suggesting it is likely to appear again.
\end{warnbox}

%% ==========================================
\section*{Bonus Cheat-Sheet --- Likely New Question Types from Lc\#11A/11B}
%% ==========================================
\begin{keybox}{If a fresh Q8-style question appears in 2026 (insurance, not yet seen in past papers)}
\begin{itemize}[noitemsep]
  \item \textbf{Draw/explain the Merkle tree} for $N$ transactions: pairwise-hash leaves up to a single root stored in the block header; explain how it gives O(1)-ish tamper detection.
  \item \textbf{ETH vs Bitcoin:} Bitcoin = purely a digital-currency/payment network (PoW, $\sim$10 min blocks); Ethereum = programmable blockchain supporting \textbf{smart contracts} (self-executing code), its currency is \textbf{ether (ETH)}, used to pay "gas" fees for computation, not just transfers.
  \item \textbf{Genesis block:} the very first block in a chain, has no predecessor (previous hash = 0000...).
  \item \textbf{Public vs Private vs Hybrid blockchain} --- table above (permission model + examples).
  \item \textbf{Blockchain vs shared database} --- table above (operations/replication/consensus/validation/confidentiality).
  \item \textbf{Limitations of blockchain (Lc\#11A p12):} higher cost than centralized DB, slower transaction throughput, smaller ledger size limit, risk of 51\% attack, irreversibility (no "undo").
  \item \textbf{Bitcoin mining mechanics:} block reward halving (every $\sim$210{,}000 blocks $\approx$4 years), self-adjusting difficulty to keep $\sim$10 min/block, mining pools to reduce reward variance.
\end{itemize}
\end{keybox}

%% ==========================================
\section*{Exam Strategy Summary}
%% ==========================================

\begin{keybox}{High-Yield Checklist for Blockchain \& Bitcoin (2/5 years confirmed, 3--8.75 marks)}
\begin{enumerate}[noitemsep]
  \item \textbf{Hashing + PoW = the core security story.} Memorize the 2024 answer pattern: chain-of-hashes makes tampering \emph{detectable}, PoW makes tampering \emph{computationally infeasible to hide} (must re-mine all subsequent blocks faster than the honest network).
  \item \textbf{Blockchain $\leftrightarrow$ cryptography (2021):} hash functions $\to$ chain integrity + Merkle tree; public-key crypto $\to$ digital signatures for ownership/authentication.
  \item \textbf{Bitcoin definition (2021):} decentralized P2P digital currency, no central authority, mining-based consensus, NOT the same thing as "blockchain."
  \item \textbf{Input/output scripts (2021):} output script = locking condition (spend requires valid signature); input script = unlocking proof (signature + pubkey); validation = combined script executes to TRUE.
  \item \textbf{Bonus prep (new Lc\#11A/11B material, not yet in past papers):} Merkle tree diagram, block-header fields (6: version/prev-hash/Merkle root/time/nBits/nonce), public vs private vs hybrid chains, blockchain vs shared DB table, ETH vs BTC, blockchain limitations.
  \item \textbf{If asked to draw something:} default to the 3-block hash-chain diagram (Block $n-1$ hash $\to$ Block $n$'s "previous hash" field) --- it's the single most reusable diagram across all sub-questions.
\end{enumerate}
\end{keybox}

\end{document}
